% Reunión focal desde I/sections/aritmetica_historias.tex, leído íntegro.
% Procedencia: technical/PROCEDENCIA_ARITMETICA_HISTORIAS_VI.md.
% Se conservan hist:seccion e hist:doble-varianza para el núcleo común.
% No incluye la representación de Weyl ni la prolongación excepcional de I.
\section{Aritmética de historias, composición y refinamiento}
\label{vi:hist:seccion-principal}

La persistencia de una partícula y la evaluación de su carácter requieren
distinguir la historia transportada de las coordenadas que se leen en ella.
El núcleo APP--TRIT--TPK proporciona los estados admisibles, sus rutas y
las operaciones que conservan su procedencia. Esta sección precisa tres
estructuras utilizadas al pasar de esos datos al atlas: la compatibilidad
entre profundidades, la composición de historias con memoria y el
transporte de observables en sentido contrario al refinamiento.
Los coeficientes algebraicos actúan sobre estados y rutas ya especificados;
las relaciones que se prueban a continuación conservan ese orden.

\subsection{Secciones compatibles y evaluación de sus lectores}
\label{hist:seccion}

Sea \(X_n\) el conjunto de estados admisibles de profundidad \(n\)
producidos por el núcleo. Un estado registra las posiciones de APP, las
hojas aditiva y multiplicativa, sus residuos y cocientes, el régimen
trítico, la orientación, la fase, la frontera, la ruta, el acarreo y la
memoria pertinente a la supervivencia. Estos datos satisfacen las
relaciones de las operaciones que los produjeron; no se consideran un
producto de coordenadas independientes.

Para \(m\ge n\), la restricción
\(r_{m,n}:X_m\longrightarrow X_n\) conserva el antecedente de profundidad
\(n\). Las restricciones satisfacen
\begin{equation}
 r_{n,n}=\operatorname{id}_{X_n},\qquad
 r_{\ell,n}=r_{m,n}\circ r_{\ell,m}
 \quad(\ell\ge m\ge n).
 \label{vi:hist:restricciones}
\end{equation}
La memoria nueva pertenece a la prolongación; la restricción conserva la
que ya estaba determinada en el prefijo.

\begin{definition}[Sección numérica compatible]
Una sección numérica HMT es una sucesión de estados de la misma genealogía
\begin{equation}
 x=(x_n)_{n\ge0}\in X_\infty:=\varprojlim_n X_n,
 \qquad r_{m,n}(x_m)=x_n\quad(m\ge n).
 \label{vi:hist:limite-inverso}
\end{equation}
Un lector es una operación definida sobre un dominio determinado de estas
secciones. Una observación finita pertenece a un prefijo; un valor escalar
es la evaluación de un lector. La sección, y la presentación de esa sección
mediante un lector, son por tanto objetos de tipos distintos.
\end{definition}

La notación de límite inverso reúne la compatibilidad de los estados ya
construidos. La existencia de una prolongación concreta se establece
mediante las reglas de admisibilidad de esa construcción, no mediante la
sola escritura de \(X_\infty\).

En el dominio arquimediano de un lector, cada prefijo determina un
intervalo cerrado no vacío \(I_n(x)=[a_n(x),b_n(x)]\), con
\begin{equation}
 I_{n+1}(x)\subseteq I_n(x),\qquad
 b_n(x)-a_n(x)\longrightarrow0.
 \label{vi:hist:cilindros}
\end{equation}
La evaluación regional por intervalos y la decisión de sus prefijos se
han reunido en la proposición~\ref{vi:dep:prefijos}.

\begin{proposition}[Valor determinado por cilindros compatibles]
Para toda sección de ese dominio existe un único número real
\(\operatorname{ev}(x)\) tal que
\begin{equation}
 \{\operatorname{ev}(x)\}=\bigcap_{n\ge0}I_n(x).
 \label{vi:hist:evaluacion}
\end{equation}
\end{proposition}
\begin{proof}
El encajamiento hace creciente la sucesión \(a_n(x)\) y decreciente
\(b_n(x)\). Ambas quedan acotadas en el intervalo inicial. Por
completitud de la recta, \(a=\sup_n a_n(x)\) y
\(b=\inf_n b_n(x)\) existen y satisfacen \(a\le b\).
La anchura tendente a cero fuerza \(a=b\). Este punto pertenece a cada
intervalo y cualquier otro punto de la intersección tendría distancia a
él menor o igual que todas las anchuras, por lo que coincidiría con él.
\end{proof}

Esta proposición determina el valor de cada sección del dominio indicado.
Una cobertura de toda \(\mathbb R\) requiere probar que cada valor
pertenece a la imagen del lector; la definición de límite inverso no
aporta por sí sola esa prueba. Del mismo modo, la identificación de dos
historias a partir de su valor requiere un resultado de inyectividad
para el lector considerado. La evaluación conserva la genealogía como
antecedente y especifica exactamente qué parte de ella publica.

\subsection{Álgebra sectorial de rutas y multiplicadores de memoria}
\label{vi:hist:algebra}

Fijemos un sector de estados y la categoría \(\mathcal C\) de sus
historias admisibles. Escribimos \(vu=v\circ u\): primero se recorre
\(u\), después \(v\). Los pasos elementales son operaciones de
selección, transporte y actualización del TPK. La categoría conserva
sus extremos y las condiciones que permiten concatenar dos rutas.

La fibra trítica local puede representarse por
\[
 \mathbb R[J_x]/(J_x^2+\tau(x)1_x),
 \qquad \tau(x)\in\{+1,0,-1\}.
\]
La representación conserva el régimen y su orientación. Para tratar
coeficientes de mayor riqueza, admitiremos un álgebra real asociativa y
unitaria \(A_x\) en cada objeto \(x\). Una transición entre regímenes
debe precisar su mapa sobre estas álgebras; no equivale a identificarlas.

La siguiente construcción se aplica a los sectores que poseen transportes
unitales \(\rho_u:A_x\to A_y\), para \(u:x\to y\), con
\begin{equation}
 \rho_{\operatorname{id}_x}=\operatorname{id}_{A_x},\qquad
 \rho_v\circ\rho_u=\rho_{vu}.
 \label{vi:hist:accion}
\end{equation}
Para cada pareja componible \(u:x\to y\), \(v:y\to z\), sea
\(\Omega(v,u)\in Z(A_z)^\times\) un multiplicador central invertible.
Sus reglas son
\begin{equation}
 \Omega(\operatorname{id}_y,u)
 =\Omega(u,\operatorname{id}_x)=1_y
 \label{vi:hist:normalizacion}
\end{equation}
y, para \(w:z\to t\),
\begin{equation}
 \rho_w\!\left(\Omega(v,u)\right)\Omega(w,vu)
 =\Omega(w,v)\Omega(wv,u).
 \label{vi:hist:cociclo}
\end{equation}
Estas condiciones describen la estructura de partida del sector: cada
aplicación concreta debe proporcionar sus transportes y multiplicadores.
Una realización mediante bimódulos conserva sus propias composiciones y
no queda reemplazada automáticamente por la acción estricta
\eqref{vi:hist:accion}.

Sea \(t(u)\) el extremo terminal de la flecha \(u\). Sobre las sumas
de soporte finito
\[
 \mathcal H=\bigoplus_{u\in\operatorname{Mor}(\mathcal C)}A_{t(u)}T_u
\]
definimos por bilinealidad
\begin{equation}
 (aT_v)(bT_u)=
 a\rho_v(b)\Omega(v,u)T_{vu}
 \label{vi:hist:producto}
\end{equation}
cuando \(u,v\) son componibles, y cero en otro caso. La suma reúne
historias con sus coeficientes; el producto concatena las historias y
conserva el multiplicador de memoria. El soporte finito asegura que cada
producto está definido sin exigir una convergencia adicional.

\subsection{Asociatividad, identidades locales y memoria del calendario}
\label{vi:hist:asociatividad}

\begin{proposition}[Composición asociativa del sector]
Bajo \eqref{vi:hist:accion}--\eqref{vi:hist:cociclo}, el producto
\eqref{vi:hist:producto} es asociativo. Si hay finitos objetos, su unidad
es \(\sum_x1_xT_{\operatorname{id}_x}\). Para un conjunto de objetos
arbitrario, las mismas expresiones sobre conjuntos finitos de extremos
dan identidades locales para toda familia finita de historias.
\end{proposition}
\begin{proof}
Sean \(u:x\to y\), \(v:y\to z\), \(w:z\to t\), con coeficientes
\(a\in A_y\), \(b\in A_z\), \(c\in A_t\). Las dos asociaciones
producen como coeficientes de \(T_{wvu}\)
\begin{align*}
 ((cT_w)(bT_v))(aT_u):\quad&
 c\rho_w(b)\Omega(w,v)\rho_{wv}(a)\Omega(wv,u),\\
 (cT_w)((bT_v)(aT_u)):\quad&
 c\rho_w(b)\rho_w\rho_v(a)
 \rho_w(\Omega(v,u))\Omega(w,vu).
\end{align*}
La acción estricta identifica \(\rho_w\rho_v(a)\) con
\(\rho_{wv}(a)\). La centralidad permite trasladar
\(\Omega(w,v)\) a la derecha de este factor, y la identidad de cociclo
iguala los factores restantes. No se permutan los coeficientes generales
\(a,b,c\), que pueden no conmutar. Si las tres flechas no forman una
cadena componible, ambas asociaciones son cero por sus extremos.
La bilinealidad extiende la igualdad a las sumas finitas. Finalmente,
\eqref{vi:hist:normalizacion} y la unitalidad de cada \(\rho_u\)
verifican la unidad. Para una familia finita basta sumar las identidades
de todos sus extremos iniciales y terminales, lo que demuestra la última
afirmación.
\end{proof}

El calendario nonádico proporciona una realización efectiva de estas
condiciones. Tomemos la categoría de un objeto con flechas
\(C_9=\mathbb Z/9\mathbb Z\), coeficientes
\(A=\mathbb R[z,z^{-1}]\) y acción identidad. Para representantes
\(0\le a,b<9\), definimos
\begin{equation}
 c_0(a,b)=\left\lfloor\frac{a+b}{9}\right\rfloor,
 \qquad\Omega(b,a)=z^{c_0(a,b)}.
 \label{vi:hist:carry-calendario}
\end{equation}
La variable formal \(z\) conserva el número de vueltas. La normalización
se sigue de \(c_0(0,a)=c_0(a,0)=0\). Si \([a+b]_9\) denota el
representante reducido, la división euclídea da
\[
 c_0(a,b)+c_0([a+b]_9,c)
 =\left\lfloor\frac{a+b+c}{9}\right\rfloor
 =c_0(b,c)+c_0(a,[b+c]_9).
\]
Esta igualdad prueba el cociclo \eqref{vi:hist:cociclo} para los
multiplicadores, y por ello la asociatividad del calendario con memoria.

Si \(T=T_{[1]_9}\), las reglas de producto dan
\(T^j=T_{[j]_9}\) para \(0\le j<9\) y \(T^9=zT_{[0]_9}\).
Al imponer posteriormente \(z^{12}=1\), resulta \(T^{108}=1\).
Las expresiones \(z^kT_{[j]_9}\), \(0\le k<12\), \(0\le j<9\),
forman las \(108\) posiciones del calendario resultante. Así se recupera
el grupo cíclico \(C_{108}\) y su proyección a \(C_9\); conservar
\(z\) sin esa relación distingue vueltas sucesivas. Ésta es una
realización de la memoria de calendario: las restantes coordenadas de
las historias TPK permanecen en su dominio propio.

\subsection{Doble variancia del refinamiento y conservación de la unidad}
\label{hist:doble-varianza}

La restricción de estados induce una prolongación de observables en sentido
contrario. Para las álgebras de funciones con producto puntual
\(D_n=\operatorname{Fun}(X_n,\mathbb R)\), la precomposición define
\begin{equation}
 r_{m,n}^*:D_n\longrightarrow D_m,
 \qquad r_{m,n}^*f=f\circ r_{m,n}.
 \label{vi:hist:precomposicion}
\end{equation}
Se conserva el orden contravariante
\(r_{\ell,n}^*=r_{\ell,m}^*\circ r_{m,n}^*\). Estas aplicaciones
son homomorfismos unitales. Si \(r_{m,n}\) es sobreyectiva, también
\(r_{m,n}^*\) es inyectiva: una diferencia entre dos funciones puede
evaluarse en una preimagen de un punto donde difieren.

\begin{lemma}[Indicadores, unidad y evaluación compatible]
Para \(x\in X_n\), sea \(e_x\) su función indicadora. Entonces
\begin{equation}
 r_{m,n}^*e_x=\mathbf1_{r_{m,n}^{-1}(x)},\qquad
 r_{m,n}^*1_n=1_m.
 \label{vi:hist:unidad}
\end{equation}
Cuando la fibra \(r_{m,n}^{-1}(x)\) es finita, la primera igualdad es
la suma finita
\[
 r_{m,n}^*e_x=\sum_{y:\,r_{m,n}(y)=x}e_y.
\]
Para toda sección \(x=(x_n)\in X_\infty\), la regla
\([f]\mapsto f(x_n)\), con \(f\in D_n\), define un homomorfismo
unital
\begin{equation}
 \operatorname{Ev}_x:D_{\mathrm{cyl}}\longrightarrow\mathbb R,
 \qquad
 D_{\mathrm{cyl}}=\varinjlim_n(D_n,r_{n+1,n}^*).
 \label{vi:hist:evaluacion-observables}
\end{equation}
\end{lemma}
\begin{proof}
Para \(y\in X_m\), se tiene
\((r_{m,n}^*e_x)(y)=1\) exactamente cuando \(r_{m,n}(y)=x\).
Ésta es la primera igualdad de \eqref{vi:hist:unidad}. Si la fibra es
finita, sus indicadores puntuales son ortogonales y suman su indicador.
La segunda igualdad se obtiene de \(1_n(r_{m,n}(y))=1\), sin ninguna
hipótesis sobre el cardinal de la fibra.

Para toda \(f\in D_n\), la naturalidad de la evaluación es
\[
 (r_{m,n}^*f)(y)=f(r_{m,n}(y)).
\]
En una sección compatible, \(r_{m,n}(x_m)=x_n\), luego
\((r_{m,n}^*f)(x_m)=f(x_n)\). Por tanto, dos representantes que se
identifican a una profundidad común tienen la misma evaluación.
La evaluación puntual preserva suma, producto y unidad, lo que concluye
la prueba de \eqref{vi:hist:evaluacion-observables}.
\end{proof}

Para fibras infinitas sigue siendo válida la formulación mediante
\(\mathbf1_{r_{m,n}^{-1}(x)}\); no se necesita interpretar una suma
infinita como una suma algebraica finita. El límite directo reúne los
observables de profundidad finita con las identificaciones inducidas por
las restricciones. La prolongación de operadores de ruta exige, además,
los cuadrados de compatibilidad de sus transportes con esas restricciones.

La unidad conservada es la función constante uno sobre el dominio de
posibilidades. Una sección individual selecciona una historia compatible
de ese dominio. El refinamiento puede aumentar la profundidad y distinguir
más prolongaciones mientras preserva la unidad y la evaluación de cada
antecedente. Esta conservación unital proporciona el contenido algebraico
preciso utilizado aquí; una medida de probabilidad, cuando interviene en
otro lector, requiere además su normalización y sus propias reglas.
